\documentclass[11pt,a4paper]{article}
\usepackage[T1]{fontenc}
\usepackage{lmodern}
\usepackage[margin=25mm]{geometry}
\usepackage{amsmath,amssymb,amsthm,mathtools}
\usepackage[hidelinks]{hyperref}
\usepackage{microtype}
\newcommand{\N}{\mathbb N}
\newcommand{\R}{\mathbb R}
\title{A compact subshift without a harmonic entropy convergence bound}
\author{Submission for TLMC problem 00000008488}
\date{24 October 2026}
\begin{document}
\maketitle

\paragraph{Claim and scope.}
The conjecture includes a uniform $1/n$ convergence rate for finite-scale
values in the equivalent definitions of Bowen topological entropy. We use
the standard convention that $n$ is the orbit length and a finite-scale value
is $n^{-1}\log r_n(U)$, where $r_n(U)$ is the minimum size of a spanning cover
at a fixed entourage $U$. We give a compact continuous dynamical system for
which even one fixed entourage admits no eventual $C/n$ error bound, for any
constant $C$. This refutes that rate assertion under this explicit conventional
interpretation. We do not dispute the classical equality of entropy definitions
on compact systems, and make no claim about the unspecified noncompact
boundary correction or periodic optimality clauses.

\paragraph{The actual compact system.}
Give $\{0,1\}^{\N}$ its product topology and usual compatible metric
\[
d(x,y)=\begin{cases}
0,&x=y,\\
2^{-\min\{j:x_j\ne y_j\}},&x\ne y.
\end{cases}
\]
Let
\[
X=\{x\in\{0,1\}^{\N}:\text{at most one coordinate of }x\text{ is }1\},
\qquad (Tx)_j=x_{j+1}.
\]
The condition defining $X$ is the intersection, over distinct $i,j$, of the
closed conditions $x_i=0$ or $x_j=0$. Hence $X$ is a closed subset of a compact
product and is compact. The left shift preserves $X$ and is continuous,
as each output coordinate is an input coordinate. Write $z$ for the all-zero
sequence and $e_j$ for the sequence whose unique $1$ occurs at $j$.

For $m\geq0$ let $U_m$ mean agreement on the first $m$ coordinates. These
are entourages, and every entourage contains some $U_m$. Indeed, agreement
on the first $m$ coordinates is equivalent to $d(x,y)\leq2^{-m}$, and
$2^{-m}$ tends to zero. Thus this family is cofinal in the actual metric
uniformity of $X$.

\paragraph{Exact spanning counts at a fixed scale.}
For an entourage $U$ set
\[
U^{(n)}=\{(x,y):(T^k x,T^k y)\in U\text{ for }0\leq k<n\}.
\]
The number $r_n(U)$ is the minimum cardinality of a finite set of centres
whose $U^{(n)}$-balls cover $X$. These are precisely the objects
\texttt{Dynamics.dynEntourage} and \texttt{Dynamics.coverMincard} used in Lean.
The actual iterates satisfy $(T^k x)_i=x_{i+k}$, so
\[
U_1^{(n)}=U_n.
\]
The $n+1$ points
\[
S_n=\{z,e_0,\ldots,e_{n-1}\}
\]
give a spanning cover: a point with its $1$ before coordinate $n$ is represented
by that $e_j$, and every other point is represented by $z$. Their $U_n$-balls
are pairwise disjoint. Consequently every spanning cover needs at least
$n+1$ centres, and
\[
r_n(U_1)=n+1.
\]
The formal proof constructs $S_n$ from the injective map
$\operatorname{Option}(\operatorname{Fin}n)\to X$ and proves its cardinality;
the covering number is not assigned as a definition. It also proves this
covering number is finite before its natural-number value is used.

\paragraph{The full topological entropy is zero.}
For any fixed $m$, the same construction with $S_{m+n}$ spans at scale
$U_m$ for $n$ iterates. Thus
\[
1\leq r_n(U_m)\leq m+n+1,
\qquad
\lim_{n\to\infty}\frac{\log(m+n+1)}{n}=0.
\]
It follows that the entropy at every $U_m$ is zero. Since the $U_m$ are
cofinal, the supremum over all entourages is also zero:
\[
h_{\rm top}(T)=0.
\]
Lean proves this equality for the native Bowen--Dinaburg quantity
\texttt{Dynamics.coverEntropy}; it is not merely a limit of an independently
chosen numerical sequence. The standard symbolic entropy convention is
also recalled in \cite[p.~1]{mitchell}, but no external entropy evaluation
is assumed in the formal proof.

\paragraph{Failure of every harmonic bound.}
For $n\geq1$, the actual finite-time value at $U_1$ is
\[
a_n=\frac{\log r_n(U_1)}{n}=\frac{\log(n+1)}n.
\]
Although $a_n\to h_{\rm top}(T)=0$, we have $na_n=\log(n+1)\to\infty$.
More explicitly, for every $C\in\R$ and $N\in\N$ choose
$n\geq\max\{N,1\}$ with $\log(n+1)>C$. Then
\[
\left|a_n-h_{\rm top}(T)\right|=\frac{\log(n+1)}n>\frac Cn.
\]
This is the full quantified obstruction to an eventual $O(1/n)$ bound,
even at this single fixed scale. A bound uniform over scales or across the
stated definitions therefore cannot hold.

\paragraph{Formal files.}
The complete construction and proof are in \texttt{lean/Main.lean}; the
conclusion is \texttt{TLMC8488.counterexample}. The project pins Lean 4.33.0
and Mathlib commit
\begin{center}
\small\texttt{db584cd6d46c92f209a44c0f1c829460d327499d}.
\end{center}
The full original bilingual source is preserved in \texttt{SOURCE.md}.
Build commands, exhaustive namespace axiom checks, fresh-kernel receipts
and file hashes accompany the submission.

\begin{thebibliography}{1}
\bibitem{mitchell}
A.~Mitchell, \emph{On word complexity and topological entropy of random
substitution subshifts}, arXiv:2305.04817v2 (2024).
\href{https://arxiv.org/abs/2305.04817}{arxiv.org/abs/2305.04817}.
\end{thebibliography}
\end{document}
